\documentclass[12pt,a4paper]{report}

% ---------------------------------------------------------
% PAQUETES
% ---------------------------------------------------------
\usepackage[utf8]{inputenc}
\usepackage[T1]{fontenc}
\usepackage[spanish,es-noquoting,es-tabla]{babel}
\usepackage{geometry}
\usepackage{setspace}
\usepackage{graphicx}
\graphicspath{{./}{imagenes/}{img/}}
\usepackage[table]{xcolor}
\usepackage{hyperref}
\usepackage{booktabs}
\usepackage{float}
\usepackage{array}
\usepackage{tabularx}
\usepackage{longtable}
\usepackage{enumitem}
\usepackage{tikz}
\usetikzlibrary{shapes.geometric,arrows.meta,positioning,calc,fit,backgrounds}

\geometry{margin=2.5cm}
\setstretch{1.5}

\hypersetup{
    colorlinks=true,
    linkcolor=blue,
    urlcolor=blue,
    citecolor=blue
}

\renewcommand{\thesection}{\arabic{section}}
\renewcommand{\thesubsection}{\thesection.\arabic{subsection}}
\renewcommand{\thesubsubsection}{\thesubsection.\arabic{subsubsection}}

% ---------------------------------------------------------
% COLORES
% ---------------------------------------------------------
\definecolor{utbazul}{HTML}{1D4ED8}
\definecolor{utbnaranja}{HTML}{EA7C1A}
\definecolor{utbverde}{HTML}{16A34A}
\definecolor{cajafondo}{HTML}{F5F6F8}
\definecolor{cajaborde}{HTML}{D9DCE3}
\definecolor{filacabecera}{HTML}{E8EDFB}
\definecolor{burbujabot}{HTML}{FFFFFF}
\definecolor{burbujauser}{HTML}{DCF8C6}
\definecolor{fondochat}{HTML}{ECE5DD}
\definecolor{grismedio}{HTML}{6B7280}

% Insignia de color con texto blanco
\newcommand{\insignia}[2]{%
  \colorbox{#1}{\textcolor{white}{\textbf{\,#2\,}}}%
}

% Caja de nota / observación
\newcommand{\nota}[1]{%
\begin{center}
\fcolorbox{cajaborde}{cajafondo}{%
\begin{minipage}{0.92\linewidth}
\vspace{0.15cm}
#1
\vspace{0.15cm}
\end{minipage}}
\end{center}}

% Estilo compacto para tablas largas
\newcommand{\tablacompacta}{\setstretch{1.1}\small}

% Captura de pantalla del prototipo. Busca el archivo en la raiz del proyecto,
% en imagenes/ y en img/; si no esta, deja un recuadro en lugar de romper la compilacion.
\newcommand{\capturaruta}{}
\newcommand{\captura}[2][width=\linewidth]{%
  \def\capturaruta{}%
  \IfFileExists{#2}{\def\capturaruta{#2}}{%
    \IfFileExists{imagenes/#2}{\def\capturaruta{imagenes/#2}}{%
      \IfFileExists{img/#2}{\def\capturaruta{img/#2}}{}}}%
  \ifx\capturaruta\empty
    \fcolorbox{cajaborde}{cajafondo}{\parbox[c][3.2cm][c]{0.85\linewidth}%
      {\centering\small\itshape Falta subir la imagen: \mbox{#2}}}%
  \else
    \fcolorbox{cajaborde}{white}{\includegraphics[#1]{\capturaruta}}%
  \fi
}

% Encabezado de tabla
\newcommand{\cabecera}[1]{\rowcolor{filacabecera}#1}

% ---------------------------------------------------------
% ESTILOS TikZ
% ---------------------------------------------------------
\tikzset{
  proc/.style={rectangle, rounded corners=3pt, draw=utbazul, fill=utbazul!8,
               text width=2.6cm, align=center, minimum height=1.0cm, font=\scriptsize},
  procm/.style={rectangle, rounded corners=3pt, draw=utbnaranja, fill=utbnaranja!10,
               text width=2.6cm, align=center, minimum height=1.0cm, font=\scriptsize},
  procs/.style={rectangle, rounded corners=3pt, draw=utbverde, fill=utbverde!10,
               text width=2.6cm, align=center, minimum height=1.0cm, font=\scriptsize},
  dec/.style={diamond, aspect=1.8, draw=grismedio, fill=gray!8,
              align=center, font=\tiny, inner sep=1pt},
  uc/.style={ellipse, draw=utbazul, fill=utbazul!7, align=center,
             font=\tiny, minimum width=2.9cm, minimum height=0.95cm, inner sep=1pt},
  flecha/.style={-{Latex[length=2mm]}, draw=grismedio},
  flechad/.style={-{Latex[length=2mm]}, dashed, draw=grismedio},
  bot/.style={rectangle, rounded corners=2pt, draw=cajaborde, fill=burbujabot,
              text width=3.6cm, align=left, font=\tiny, inner sep=3pt},
  usr/.style={rectangle, rounded corners=2pt, draw=cajaborde, fill=burbujauser,
              text width=2.6cm, align=left, font=\tiny, inner sep=3pt},
}

% Actor (monigote) para el diagrama de casos de uso
\newcommand{\actor}[3]{%
  \begin{scope}[shift={(#1,#2)}, scale=0.30]
    \draw[thick] (0,0.55) circle (0.32);
    \draw[thick] (0,0.23) -- (0,-0.7);
    \draw[thick] (-0.55,-0.15) -- (0.55,-0.15);
    \draw[thick] (0,-0.7) -- (-0.42,-1.5);
    \draw[thick] (0,-0.7) -- (0.42,-1.5);
  \end{scope}
  \node[font=\tiny, align=center] at (#1,{#2-0.72}) {#3};
}

\begin{document}

% ---------------------------------------------------------
% PORTADA
% ---------------------------------------------------------
\begin{titlepage}
    \centering

    \begin{minipage}{0.45\textwidth}
        \centering
        \IfFileExists{logo_utb.png}{\includegraphics[width=0.6\textwidth]{logo_utb.png}}{\vspace{2.2cm}}
    \end{minipage}
    \hfill
    \begin{minipage}{0.45\textwidth}
        \centering
        \IfFileExists{logo_escuela.png}{\includegraphics[width=0.8\textwidth]{logo_escuela.png}}{\vspace{2.2cm}}
    \end{minipage}

    \vspace{1.5cm}

    {\Large Universidad Tecnológica de Bolívar\\[0.5cm]
    Facultad de Ingeniería}\\[2cm]

    {\bfseries\LARGE Tooli: Asistente Conversacional de WhatsApp para la Captación de Aspirantes a Posgrado\\[0.5cm]}
    {\large Entregable I: Diccionario de datos, mapa de procesos, especificación de requisitos (IEEE 830), historias de usuario, casos de uso y mockup}\\[2cm]

    {\large Estudiantes:\\
    \textbf{Manuel Esteban Severiche Puello}\\
    \textbf{Jerson David Díaz Álvarez}\\
    \textbf{Ángel Jhoan Acero Martínez}}\\[0.8cm]

    {\large Curso:\\
    \textbf{Proyecto de Ingeniería 2}}\\[0.8cm]

    {\large Docente:\\
    \textbf{Malorys Mercedes Elles Fang}}\\[0.8cm]

    \vfill

    {\large Cartagena de Indias, \today}
\end{titlepage}

% ---------------------------------------------------------
% ÍNDICE
% ---------------------------------------------------------
\tableofcontents
\clearpage

% =========================================================
% 1. INTRODUCCIÓN
% =========================================================
\section{Introducción}

\subsection{Contexto}

La Universidad Tecnológica de Bolívar (UTB) ofrece un portafolio amplio de programas de posgrado
(especializaciones, maestrías y doctorados) cuya captación de aspirantes depende hoy de canales
con alta fricción: formularios web, correo electrónico y llamadas telefónicas. El aspirante que
descubre un programa en una feria, en una publicación en redes o en un evento institucional debe
salir de su contexto, abrir un portal, diligenciar un formulario y esperar a que alguien lo
contacte. Ese salto de canal es donde se pierde la mayoría de los interesados.

En Colombia, WhatsApp es el canal de comunicación cotidiano por excelencia. Atender por allí elimina
la fricción de instalar una aplicación o navegar un portal: el aspirante escribe donde ya está.

\subsection{Problema}

La Coordinación de Admisiones de Posgrados enfrenta tres problemas concretos:

\begin{enumerate}[itemsep=2pt]
    \item \textbf{Pérdida de interesados por fricción de canal.} Entre el momento en que una persona
    se interesa y el momento en que un asesor la contacta pueden pasar horas o días; buena parte de
    esos contactos nunca se concreta.
    \item \textbf{Repetición de preguntas de bajo valor.} Costos, duración, modalidad, requisitos y
    fechas se responden manualmente una y otra vez, consumiendo el tiempo de los asesores.
    \item \textbf{Falta de trazabilidad del prospecto.} No queda registro estructurado de quién
    mostró interés, en qué programa y en qué momento, lo que impide medir la conversión y hacer
    seguimiento.
\end{enumerate}

\subsection{Solución propuesta}

\textbf{Tooli} es un asistente conversacional sobre WhatsApp que atiende al aspirante en el mismo
canal donde surge su interés: le muestra el catálogo de posgrados con costos y duración reales,
resuelve dudas abiertas mediante un asistente de inteligencia artificial apoyado en una base de
conocimiento institucional y, cuando el interés se concreta, recoge nombre, correo y programa
de interés y entrega el prospecto ya calificado a un asesor humano en la plataforma de atención,
sincronizándolo además con el CRM de la universidad.

\subsection{Objetivos}

\paragraph{Objetivo general.} Desarrollar un asistente conversacional de WhatsApp que automatice la
información sobre programas de posgrado de la UTB y la captación y entrega trazable de prospectos a
la Coordinación de Admisiones.

\paragraph{Objetivos específicos.}
\begin{enumerate}[itemsep=2pt]
    \item Modelar los elementos de información del dominio de captación de posgrados en un
    diccionario de datos.
    \item Levantar el mapa de procesos del negocio de captación y ubicar en él el alcance del
    sistema.
    \item Especificar los requisitos funcionales y no funcionales bajo el estándar IEEE 830.
    \item Expresar las necesidades de los usuarios como historias de usuario con criterios de
    aceptación verificables.
    \item Modelar el comportamiento del sistema mediante casos de uso.
    \item Prototipar la interfaz conversacional y el tablero de asesores mediante mockups.
\end{enumerate}

\subsection{Alcance de este entregable}

Este documento corresponde al \textbf{Entregable I} del curso y cubre exclusivamente la fase de
análisis y especificación del \textbf{flujo de posgrados} de Tooli. Quedan expresamente fuera de
este documento los demás módulos que comparten la misma base técnica: los servicios al estudiante
(notas, turno de matrícula y recibo de pago) y otros bots alojados en el mismo servidor, que serán
objeto de documentación independiente.

\clearpage

% =========================================================
% 2. DICCIONARIO DE DATOS
% =========================================================
\section{Diccionario de Datos (Elementos de Información)}

El diccionario de datos describe los elementos de información que el sistema captura, transforma,
almacena y transmite en el flujo de posgrados. Se organiza por entidad, indicando para cada atributo
su tipo, formato, obligatoriedad y regla de negocio asociada.

\subsection{Convenciones}

\begin{table}[H]
    \centering
    \tablacompacta
    \begin{tabular}{ll}
        \toprule
        \textbf{Marca} & \textbf{Significado} \\
        \midrule
        \textbf{PK} & Identificador único de la entidad. \\
        \textbf{FK} & Referencia a otra entidad. \\
        Sí / No & Atributo obligatorio o no obligatorio. \\
        \bottomrule
    \end{tabular}
\end{table}

\subsection{Entidad: Prospecto}

Persona interesada en un programa de posgrado que completa la captura de datos. Es el elemento de
información central del sistema.

{\tablacompacta
\begin{longtable}{p{3.1cm}p{2.2cm}p{1.3cm}p{6.0cm}}
\caption{Diccionario de datos: Prospecto.}\label{tab:dd-prospecto}\\
\toprule
\cabecera{\textbf{Atributo}} & \textbf{Tipo} & \textbf{Oblig.} & \textbf{Descripción y regla de negocio} \\
\midrule
\endfirsthead
\toprule
\textbf{Atributo} & \textbf{Tipo} & \textbf{Oblig.} & \textbf{Descripción y regla de negocio} \\
\midrule
\endhead
\bottomrule
\endfoot
whatsapp \textbf{(PK)} & Cadena (12--15) & Sí & Número en formato E.164 sin el signo ``+'' (ej.: 573001234567). Lo entrega la plataforma de mensajería; identifica de forma única al prospecto. \\
nombre & Cadena (3--80) & Sí & Nombre completo declarado por el usuario. Mínimo 3 caracteres; se rechaza y se vuelve a solicitar si no cumple. \\
correo & Cadena (5--120) & Sí & Correo electrónico de contacto. Validado con expresión regular de formato; se normaliza a minúsculas. \\
programa\_interes & Cadena (1--120) & Sí & Programa o área declarada de interés. Admite texto libre, específico (``MBA'') o general (``algo de ingeniería'', ``no sé aún''). \\
categoria\_vista & Enumerado & No & Última categoría consultada antes de solicitar asesor: \textit{Especializaciones}, \textit{Maestrías}, \textit{Doctorados}. Se deriva de la navegación, no se pregunta. \\
interes\_financiacion & Enumerado & No & Respuesta a la pregunta de interés en financiación: \textit{Sí}, \textit{No}, \textit{Tal vez}. Aplica al flujo de captación por evento. \\
consentimiento & Booleano & Sí & Autorización explícita de tratamiento de datos (Ley 1581 de 2012). Si es falso, no se persiste ningún otro dato del prospecto. \\
fecha\_registro & Fecha y hora & Sí & Marca temporal de la captura, en zona horaria \textit{America/Bogota}. Generada por el sistema. \\
origen & Enumerado & Sí & Canal de entrada: \textit{menú del bot}, \textit{código QR del evento}, \textit{landing pública}. \\
estado & Enumerado & Sí & Ciclo de vida: \textit{capturado}, \textit{entregado a asesor}, \textit{en atención}, \textit{cerrado}. \\
\end{longtable}}

\subsection{Entidad: Programa de posgrado}

Elemento del catálogo que el bot presenta al aspirante. Es información de solo lectura para el
usuario final.

{\tablacompacta
\begin{longtable}{p{3.1cm}p{2.2cm}p{1.3cm}p{6.0cm}}
\caption{Diccionario de datos: Programa de posgrado.}\label{tab:dd-programa}\\
\toprule
\cabecera{\textbf{Atributo}} & \textbf{Tipo} & \textbf{Oblig.} & \textbf{Descripción y regla de negocio} \\
\midrule
\endfirsthead
\toprule
\textbf{Atributo} & \textbf{Tipo} & \textbf{Oblig.} & \textbf{Descripción y regla de negocio} \\
\midrule
\endhead
\bottomrule
\endfoot
codigo \textbf{(PK)} & Cadena (8) & Sí & Identificador interno del programa. \\
nombre & Cadena (5--100) & Sí & Denominación oficial del programa. \\
categoria & Enumerado & Sí & \textit{Especialización}, \textit{Maestría} o \textit{Doctorado}. Determina la sección del catálogo en la que aparece. \\
modalidad & Enumerado & Sí & \textit{Presencial} o \textit{Virtual}. \\
duracion\_semestres & Entero (1--10) & Sí & Duración del plan de estudios. \\
valor\_semestre & Moneda (COP) & Sí & Valor por semestre vigente. Los doctorados pueden expresarse en SMMLV. \\
valor\_inscripcion & Moneda (COP) & Sí & Valor del derecho de inscripción, común por categoría. \\
requisitos & Texto & No & Requisitos de admisión (ej.: título de maestría o pregrado con alto promedio para doctorado). \\
vigencia & Fecha & Sí & Período académico al que corresponden los valores publicados. \\
\end{longtable}}

\subsection{Entidad: Sesión conversacional}

Estado de la conversación de un usuario. Permite diálogos de varios turnos sin perder el hilo. Es
información volátil: vive en memoria con tiempo de vida limitado.

{\tablacompacta
\begin{longtable}{p{3.1cm}p{2.2cm}p{1.3cm}p{6.0cm}}
\caption{Diccionario de datos: Sesión conversacional.}\label{tab:dd-sesion}\\
\toprule
\cabecera{\textbf{Atributo}} & \textbf{Tipo} & \textbf{Oblig.} & \textbf{Descripción y regla de negocio} \\
\midrule
\endfirsthead
\toprule
\textbf{Atributo} & \textbf{Tipo} & \textbf{Oblig.} & \textbf{Descripción y regla de negocio} \\
\midrule
\endhead
\bottomrule
\endfoot
clave \textbf{(PK)} & Cadena & Sí & Clave compuesta \texttt{session:\{whatsapp\}}. \\
step & Cadena (5--30) & Sí & Paso actual de la máquina de estados (ej.: \texttt{menu}, \texttt{programas\_categoria}, \texttt{prospecto\_nombre}, \texttt{chatting\_with\_ai}, \texttt{with\_agent}). Determina qué manejador atiende el siguiente mensaje. \\
data & Diccionario & No & Datos parciales acumulados durante el diálogo (nombre, correo, categoría vista, identificador de conversación). \\
ttl & Entero (segundos) & Sí & Tiempo de vida restante. Al expirar, la sesión se elimina y el usuario vuelve al menú principal. \\
\end{longtable}}

\subsection{Entidad: Mensaje}

Unidad de intercambio entre el usuario y el sistema. No se persiste el contenido: se procesa y se
descarta.

{\tablacompacta
\begin{longtable}{p{3.1cm}p{2.2cm}p{1.3cm}p{6.0cm}}
\caption{Diccionario de datos: Mensaje.}\label{tab:dd-mensaje}\\
\toprule
\cabecera{\textbf{Atributo}} & \textbf{Tipo} & \textbf{Oblig.} & \textbf{Descripción y regla de negocio} \\
\midrule
\endfirsthead
\toprule
\textbf{Atributo} & \textbf{Tipo} & \textbf{Oblig.} & \textbf{Descripción y regla de negocio} \\
\midrule
\endhead
\bottomrule
\endfoot
message\_id \textbf{(PK)} & Cadena & Sí & Identificador asignado por la plataforma de mensajería. Se marca como procesado durante 5 minutos para garantizar idempotencia ante reenvíos del webhook. \\
from \textbf{(FK)} & Cadena & Sí & Número del remitente; referencia al Prospecto o al usuario anónimo. \\
direccion & Enumerado & Sí & \textit{Entrante} o \textit{Saliente}. \\
tipo & Enumerado & Sí & \textit{texto}, \textit{lista interactiva}, \textit{botones}, \textit{respuesta interactiva}. \\
contenido & Texto & Sí & Cuerpo del mensaje o identificador de la opción seleccionada. No se almacena en texto plano. \\
timestamp & Fecha y hora & Sí & Marca temporal de recepción o envío. \\
\end{longtable}}

\subsection{Entidad: Conversación de asesoría}

Registro de la atención humana abierta en la plataforma de atención (Chatwoot) cuando el prospecto
solicita hablar con un asesor.

{\tablacompacta
\begin{longtable}{p{3.1cm}p{2.2cm}p{1.3cm}p{6.0cm}}
\caption{Diccionario de datos: Conversación de asesoría.}\label{tab:dd-conversacion}\\
\toprule
\cabecera{\textbf{Atributo}} & \textbf{Tipo} & \textbf{Oblig.} & \textbf{Descripción y regla de negocio} \\
\midrule
\endfirsthead
\toprule
\textbf{Atributo} & \textbf{Tipo} & \textbf{Oblig.} & \textbf{Descripción y regla de negocio} \\
\midrule
\endhead
\bottomrule
\endfoot
conversation\_id \textbf{(PK)} & Entero & Sí & Identificador de la conversación en la plataforma de atención. \\
contact\_id \textbf{(FK)} & Entero & Sí & Contacto asociado; se crea si no existe y se reutiliza si ya existía. \\
whatsapp \textbf{(FK)} & Cadena & Sí & Número del prospecto atendido. \\
equipo & Enumerado & Sí & Equipo destinatario. En el flujo de posgrados siempre \textit{Admisiones}. \\
etiquetas & Lista & Sí & Marcas de clasificación; el flujo asigna \texttt{posgrado-prospecto}. \\
prioridad & Enumerado & Sí & \textit{baja}, \textit{media}, \textit{alta}. El flujo asigna \textit{media} por defecto. \\
mensaje\_contexto & Texto & Sí & Ficha inicial que ve el asesor: nombre, WhatsApp, correo, programa de interés, categoría vista y fecha de recepción. \\
atributos & Diccionario & No & Área, programa de interés y correo, como atributos personalizados consultables. \\
estado & Enumerado & Sí & \textit{abierta}, \textit{pendiente}, \textit{resuelta}. \\
\end{longtable}}

\subsection{Entidad: Contacto en el CRM}

Proyección del prospecto en el CRM institucional, para seguimiento comercial posterior.

{\tablacompacta
\begin{longtable}{p{3.1cm}p{2.2cm}p{1.3cm}p{6.0cm}}
\caption{Diccionario de datos: Contacto en el CRM.}\label{tab:dd-crm}\\
\toprule
\cabecera{\textbf{Atributo}} & \textbf{Tipo} & \textbf{Oblig.} & \textbf{Descripción y regla de negocio} \\
\midrule
\endfirsthead
\toprule
\textbf{Atributo} & \textbf{Tipo} & \textbf{Oblig.} & \textbf{Descripción y regla de negocio} \\
\midrule
\endhead
\bottomrule
\endfoot
email \textbf{(PK)} & Cadena & Sí & Clave de deduplicación: la sincronización es una operación de inserción o actualización por correo. \\
firstname & Cadena & Sí & Primera palabra del nombre declarado. \\
lastname & Cadena & No & Resto del nombre declarado. \\
phone & Cadena & Sí & Número en formato E.164 con el signo ``+''. \\
posgrado\_interes & Cadena & No & Programa declarado; propiedad personalizada. \\
fuente & Cadena & No & Origen del contacto (ej.: \textit{WhatsApp Tooli}); propiedad personalizada. \\
consentimiento & Booleano & Sí & Registro de la autorización de tratamiento de datos. \\
interes\_financiacion & Cadena & No & Interés declarado en opciones de financiación. \\
\end{longtable}}

\nota{\textbf{Regla de degradación:} si alguna propiedad personalizada no existe en el CRM, la
sincronización se reintenta únicamente con las propiedades estándar, de modo que el contacto nunca
se pierda por un problema de configuración.}

\subsection{Entidad: Registro en hoja de captación}

Fila del registro tabular usado en los eventos de captación, accesible para la Coordinación sin
necesidad de entrar al sistema.

{\tablacompacta
\begin{longtable}{p{2.0cm}p{3.6cm}p{2.0cm}p{5.0cm}}
\caption{Diccionario de datos: Registro en hoja de captación.}\label{tab:dd-hoja}\\
\toprule
\cabecera{\textbf{Col.}} & \textbf{Campo} & \textbf{Tipo} & \textbf{Descripción} \\
\midrule
\endfirsthead
\toprule
\textbf{Col.} & \textbf{Campo} & \textbf{Tipo} & \textbf{Descripción} \\
\midrule
\endhead
\bottomrule
\endfoot
A & Fecha y hora & Fecha y hora & Momento del registro (\textit{America/Bogota}). \\
B & WhatsApp & Cadena & Número del prospecto. \\
C & Nombre & Cadena & Nombre completo declarado. \\
D & Correo & Cadena & Correo de contacto. \\
E & Interés en financiación & Enumerado & Sí / No / Tal vez. \\
F & Posgrado de interés & Cadena & Se completa en un paso posterior del diálogo. \\
G & Autoriza contacto & Booleano & Se completa al final del diálogo. \\
\end{longtable}}

\subsection{Entidad: Asesor / Usuario del tablero}

Persona de la Coordinación de Admisiones con acceso al tablero interno.

{\tablacompacta
\begin{longtable}{p{3.1cm}p{2.2cm}p{1.3cm}p{6.0cm}}
\caption{Diccionario de datos: Asesor / Usuario del tablero.}\label{tab:dd-asesor}\\
\toprule
\cabecera{\textbf{Atributo}} & \textbf{Tipo} & \textbf{Oblig.} & \textbf{Descripción y regla de negocio} \\
\midrule
\endfirsthead
\toprule
\textbf{Atributo} & \textbf{Tipo} & \textbf{Oblig.} & \textbf{Descripción y regla de negocio} \\
\midrule
\endhead
\bottomrule
\endfoot
usuario \textbf{(PK)} & Cadena (4--30) & Sí & Nombre de acceso, único. \\
hash\_clave & Cadena & Sí & Resumen criptográfico de la contraseña. Nunca se almacena la contraseña en claro. \\
nombre & Cadena & Sí & Nombre para mostrar. \\
rol & Enumerado & Sí & \textit{asesor} o \textit{administrador}. El rol determina las secciones visibles del tablero. \\
activo & Booleano & Sí & Permite revocar el acceso sin borrar el registro. \\
ultimo\_acceso & Fecha y hora & No & Fecha del último inicio de sesión exitoso. \\
\end{longtable}}

\subsection{Entidad: Evento de métrica}

Unidad de medición del uso del sistema. No contiene datos personales.

{\tablacompacta
\begin{longtable}{p{3.1cm}p{2.2cm}p{1.3cm}p{6.0cm}}
\caption{Diccionario de datos: Evento de métrica.}\label{tab:dd-metrica}\\
\toprule
\cabecera{\textbf{Atributo}} & \textbf{Tipo} & \textbf{Oblig.} & \textbf{Descripción y regla de negocio} \\
\midrule
\endfirsthead
\toprule
\textbf{Atributo} & \textbf{Tipo} & \textbf{Oblig.} & \textbf{Descripción y regla de negocio} \\
\midrule
\endhead
\bottomrule
\endfoot
nombre \textbf{(PK)} & Cadena & Sí & Identificador del evento (ej.: \texttt{programas\_consultado}, \texttt{prospecto\_completado}, \texttt{handoff\_asesor}). \\
fecha & Fecha & Sí & Día al que se agrega el contador. \\
valor & Entero & Sí & Número de ocurrencias acumuladas. \\
\end{longtable}}

\subsection{Relaciones entre entidades}

\begin{table}[H]
    \centering
    \tablacompacta
    \begin{tabularx}{\textwidth}{llX}
        \toprule
        \cabecera{\textbf{Origen}} & \textbf{Destino} & \textbf{Relación} \\
        \midrule
        Prospecto & Sesión conversacional & 1 : 1, mientras la sesión esté vigente. \\
        Prospecto & Conversación de asesoría & 1 : N, puede solicitar asesor más de una vez. \\
        Prospecto & Contacto en el CRM & 1 : 1, deduplicado por correo. \\
        Prospecto & Registro en hoja & 1 : 1, solo en el flujo de captación por evento. \\
        Programa & Prospecto & 1 : N, un programa concentra varios interesados. \\
        Asesor & Conversación de asesoría & 1 : N, un asesor atiende varias conversaciones. \\
        Sesión & Mensaje & 1 : N, los mensajes se procesan en el contexto de la sesión. \\
        \bottomrule
    \end{tabularx}
    \caption{Relaciones entre los elementos de información.}
\end{table}

\clearpage
% =========================================================
% 3. MAPA DE PROCESOS
% =========================================================
\section{Mapa de Procesos}

El mapa de procesos ubica el sistema dentro del negocio de captación de aspirantes a posgrado de la
UTB. Se emplea la clasificación clásica en procesos estratégicos, misionales (o de cadena de valor)
y de apoyo. El sistema \textbf{Tooli} no reemplaza el proceso misional: lo automatiza en sus tres
primeras etapas y garantiza la trazabilidad de la cuarta.

\subsection{Mapa general}

\begin{figure}[H]
\centering
\begin{tikzpicture}[node distance=0.45cm and 0.45cm]

% --- Procesos estratégicos ---
\node[proc] (e1) {Planeación de la oferta de posgrados};
\node[proc, right=of e1] (e2) {Estrategia de captación y mercadeo};
\node[proc, right=of e2] (e3) {Gobierno de datos y cumplimiento legal};
\node[proc, right=of e3] (e4) {Seguimiento de indicadores de conversión};

\node[font=\scriptsize\bfseries, left=0.3cm of e1, text width=1.7cm, align=right,
      text=utbazul] (le) {PROCESOS ESTRATÉGICOS};

% --- Procesos misionales ---
\node[procm, below=1.5cm of e1] (m1) {1. Atracción y primer contacto};
\node[procm, right=of m1] (m2) {2. Información y resolución de dudas};
\node[procm, right=of m2] (m3) {3. Calificación y captura del prospecto};
\node[procm, right=of m3] (m4) {4. Asesoría humana y seguimiento};

\node[font=\scriptsize\bfseries, left=0.3cm of m1, text width=1.7cm, align=right,
      text=utbnaranja] (lm) {PROCESOS MISIONALES};

% --- Procesos de apoyo ---
\node[procs, below=1.5cm of m1] (a1) {Infraestructura y despliegue};
\node[procs, right=of a1] (a2) {Gestión de la base de conocimiento};
\node[procs, right=of a2] (a3) {Administración de asesores y accesos};
\node[procs, right=of a3] (a4) {Observabilidad y soporte};

\node[font=\scriptsize\bfseries, left=0.3cm of a1, text width=1.7cm, align=right,
      text=utbverde] (la) {PROCESOS DE APOYO};

% --- Flechas de la cadena de valor ---
\draw[flecha] (m1) -- (m2);
\draw[flecha] (m2) -- (m3);
\draw[flecha] (m3) -- (m4);

% --- Entrada y salida ---
\node[font=\scriptsize\itshape, left=0.9cm of lm, align=center, text width=1.6cm]
  (entrada) {Interés del aspirante};
\node[font=\scriptsize\itshape, right=0.5cm of m4, align=center, text width=1.6cm]
  (salida) {Aspirante inscrito};
\draw[flecha] (entrada) -- (lm);
\draw[flecha] (m4) -- (salida);

% --- Influencia vertical ---
\draw[flechad] (e1.south) -- ($(m1.north)+(0,0)$);
\draw[flechad] (e2.south) -- (m2.north);
\draw[flechad] (e3.south) -- (m3.north);
\draw[flechad] (e4.south) -- (m4.north);
\draw[flechad] (a1.north) -- (m1.south);
\draw[flechad] (a2.north) -- (m2.south);
\draw[flechad] (a3.north) -- (m3.south);
\draw[flechad] (a4.north) -- (m4.south);

\end{tikzpicture}
\caption{Mapa de procesos de la captación de aspirantes a posgrado.}
\label{fig:mapa-procesos}
\end{figure}

\subsection{Descripción de los procesos misionales}

\begin{table}[H]
    \centering
    \tablacompacta
    \begin{tabularx}{\textwidth}{p{2.6cm}XX}
        \toprule
        \cabecera{\textbf{Proceso}} & \textbf{Entradas y actividades} & \textbf{Salidas} \\
        \midrule
        1. Atracción y primer contacto &
        Campaña, feria, código QR o publicación en redes. El aspirante escribe al número de WhatsApp
        de la universidad y recibe el menú principal. &
        Conversación iniciada; sesión creada. \\[0.2cm]
        2. Información y resolución de dudas &
        Consulta del catálogo por categoría, detalle del programa con costos y duración, y preguntas
        abiertas atendidas por el asistente de IA sobre la base de conocimiento institucional. &
        Aspirante informado; categoría de interés identificada. \\[0.2cm]
        3. Calificación y captura del prospecto &
        Solicitud de asesor, consentimiento de tratamiento de datos, captura de nombre, correo y
        programa de interés, con validación de formato. &
        Prospecto calificado y registrado. \\[0.2cm]
        4. Asesoría humana y seguimiento &
        Apertura de conversación en la plataforma de atención con la ficha del prospecto,
        conversación puente entre asesor y aspirante, y sincronización con el CRM. &
        Prospecto atendido y trazable; insumo para inscripción. \\
        \bottomrule
    \end{tabularx}
    \caption{Caracterización de los procesos misionales.}
\end{table}

\subsection{Flujo detallado del proceso misional automatizado}

El siguiente diagrama detalla el recorrido del aspirante dentro del sistema, desde el mensaje
inicial hasta la entrega del prospecto al asesor.

\begin{figure}[H]
\centering
\begin{tikzpicture}[node distance=0.55cm and 0.7cm]

\node[proc, fill=gray!8, draw=grismedio] (ini) {El aspirante escribe por WhatsApp};
\node[proc, below=of ini] (menu) {Se muestra el menú principal};
\node[dec, below=of menu, text width=1.6cm] (opc) {¿Qué opción elige?};

\node[procm, left=1.5cm of opc] (cat) {Selecciona categoría de programa};
\node[procm, below=1.6cm of opc] (ia) {Pregunta al asistente de IA};
\node[procm, right=1.5cm of opc] (ases) {Solicita hablar con un asesor};

\node[procm, below=of cat] (det) {Ve el detalle: costos, duración y modalidad};

\node[dec, below=of ia, text width=1.6cm] (segue) {¿Quiere continuar?};

\node[procm, below=1.2cm of ases] (cons) {Se solicita consentimiento de datos};
\node[dec, below=of cons, text width=1.6cm] (auth) {¿Autoriza?};
\node[procm, below=of auth] (cap) {Captura de nombre, correo y programa};
\node[proc, fill=gray!8, draw=grismedio, below=of cap] (fin2) {Fin: sin registro};

\node[procm, below=2.4cm of segue] (hand) {Se abre conversación con el equipo de Admisiones};
\node[procs, below=of hand] (crm) {Se sincroniza el prospecto con el CRM};
\node[proc, fill=gray!8, draw=grismedio, below=of crm] (fin) {El asesor atiende al aspirante};

\draw[flecha] (ini) -- (menu);
\draw[flecha] (menu) -- (opc);
\draw[flecha] (opc) -- node[font=\tiny, above] {Programas} (cat);
\draw[flecha] (opc) -- node[font=\tiny, left] {IA} (ia);
\draw[flecha] (opc) -- node[font=\tiny, above] {Asesor} (ases);
\draw[flecha] (cat) -- (det);
\draw[flecha] (det) -- (segue);
\draw[flecha] (ia) -- (segue);
\draw[flecha] (segue) -- node[font=\tiny, left] {Sí, asesor} (hand);
\draw[flecha] (ases) -- (cons);
\draw[flecha] (cons) -- (auth);
\draw[flecha] (auth) -- node[font=\tiny, left] {Sí} (cap);
\draw[flecha] (auth.east) -- ++(1.1,0) |- node[font=\tiny, above, pos=0.75] {No} (fin2.east);
\draw[flecha] (cap.west) -- ++(-0.6,0) |- (hand.east);
\draw[flecha] (hand) -- (crm);
\draw[flecha] (crm) -- (fin);

\end{tikzpicture}
\caption{Flujo detallado del proceso de captación automatizado.}
\label{fig:flujo-detallado}
\end{figure}

\nota{\textbf{Alcance del sistema en el mapa:} Tooli automatiza por completo los procesos misionales
1, 2 y 3, y habilita el 4 entregando el prospecto con contexto al asesor humano. Los procesos
estratégicos permanecen bajo responsabilidad de la Coordinación de Admisiones; los procesos de apoyo
son responsabilidad del equipo técnico del proyecto.}

\clearpage

% =========================================================
% 4. ESPECIFICACIÓN DE REQUISITOS (IEEE 830)
% =========================================================
\section{Especificación de Requisitos (IEEE 830)}

\subsection{Propósito}

Este apartado especifica los requisitos funcionales y no funcionales del flujo de posgrados de
Tooli, siguiendo la estructura del estándar IEEE 830 para especificaciones de requisitos de
software. Está dirigido al equipo de desarrollo, a la docente del curso y a la Coordinación de
Admisiones de Posgrados como parte interesada principal.

\subsection{Alcance del producto}

\paragraph{Nombre del producto.} Tooli, asistente conversacional de posgrados UTB.

\paragraph{Qué hace el producto.}
\begin{itemize}[itemsep=2pt]
    \item Atiende por WhatsApp a personas interesadas en los programas de posgrado de la UTB, las
    24 horas.
    \item Presenta el catálogo de especializaciones, maestrías y doctorados con costos, duración y
    modalidad vigentes.
    \item Resuelve dudas abiertas mediante un asistente de IA apoyado en una base de conocimiento
    institucional.
    \item Captura los datos del aspirante interesado previo consentimiento informado.
    \item Entrega el prospecto calificado a un asesor humano y mantiene una conversación puente
    entre ambos.
    \item Sincroniza el prospecto con el CRM institucional y expone métricas de uso.
\end{itemize}

\paragraph{Qué no hace el producto.}
\begin{itemize}[itemsep=2pt]
    \item No formaliza inscripciones ni matrículas; no procesa pagos.
    \item No emite decisiones de admisión ni valida documentos académicos.
    \item No sustituye al asesor humano: lo antecede y lo alimenta.
    \item No cubre los servicios al estudiante activo (notas, turno, recibo), que son objeto de otro
    documento.
\end{itemize}

\paragraph{Beneficios esperados.} Reducción del tiempo entre el interés y el primer contacto humano;
liberación del tiempo de los asesores de preguntas repetitivas; trazabilidad completa del embudo de
captación.

\subsection{Perspectiva del producto}

Tooli es un producto \textbf{nuevo}, no un reemplazo de un sistema existente, pero \textbf{no es
autónomo}: opera como pieza de un ecosistema y depende de servicios externos. La figura
\ref{fig:contexto} muestra su contexto.

\begin{figure}[H]
\centering
\begin{tikzpicture}[node distance=0.9cm and 1.3cm]

\node[proc, minimum height=1.6cm, text width=3.1cm, fill=utbazul!15, very thick] (sis)
  {\textbf{Tooli}\\ Backend conversacional\\ (máquina de estados)};

\node[proc, above=1.3cm of sis, draw=grismedio, fill=gray!8] (wa) {WhatsApp Cloud API (Meta)};
\node[proc, left=of sis, draw=grismedio, fill=gray!8] (redis) {Redis\\ (sesiones y métricas)};
\node[proc, right=of sis, draw=grismedio, fill=gray!8] (ia) {Servicio de IA\\ (modelo de lenguaje)};
\node[proc, below left=1.3cm and 0.1cm of sis, draw=grismedio, fill=gray!8] (cw) {Chatwoot\\ (atención humana)};
\node[proc, below right=1.3cm and 0.1cm of sis, draw=grismedio, fill=gray!8] (crm) {CRM HubSpot};
\node[proc, below=2.9cm of sis, draw=grismedio, fill=gray!8] (sheets) {Google Sheets\\ (registro de captación)};

\node[font=\scriptsize, above=0.25cm of wa] (asp) {Aspirante};
\node[proc, left=1.3cm of cw, draw=grismedio, fill=gray!8] (dash) {Tablero de asesores\\ (web, JWT)};

\draw[flecha] (asp) -- (wa);
\draw[{Latex[length=2mm]}-{Latex[length=2mm]}, draw=grismedio] (wa) -- (sis);
\draw[{Latex[length=2mm]}-{Latex[length=2mm]}, draw=grismedio] (sis) -- (redis);
\draw[flecha] (sis) -- (ia);
\draw[{Latex[length=2mm]}-{Latex[length=2mm]}, draw=grismedio] (sis) -- (cw);
\draw[flecha] (sis) -- (crm);
\draw[flecha] (sis) -- (sheets);
\draw[{Latex[length=2mm]}-{Latex[length=2mm]}, draw=grismedio] (dash) -- (sis);

\end{tikzpicture}
\caption{Diagrama de contexto del sistema.}
\label{fig:contexto}
\end{figure}

\paragraph{Interfaces de sistema.}
\begin{itemize}[itemsep=2pt]
    \item \textbf{WhatsApp Cloud API (Meta):} recepción de mensajes por webhook y envío de mensajes
    de texto, listas y botones interactivos.
    \item \textbf{Redis:} almacenamiento del estado conversacional con tiempo de vida limitado, y de
    los contadores de métricas.
    \item \textbf{Servicio de modelo de lenguaje:} generación de respuestas del asistente de IA.
    \item \textbf{Chatwoot:} creación de contactos y conversaciones, y puente bidireccional de
    mensajes con el asesor.
    \item \textbf{HubSpot:} inserción o actualización del contacto por correo.
    \item \textbf{Google Sheets:} registro tabular de la captación en eventos.
\end{itemize}

\paragraph{Interfaces de usuario.}
\begin{itemize}[itemsep=2pt]
    \item \textbf{Interfaz conversacional} en WhatsApp: texto, listas interactivas (hasta 10 filas) y
    botones (hasta 3 por mensaje).
    \item \textbf{Tablero web de asesores:} aplicación de una sola página con autenticación por
    token, para métricas y administración.
    \item \textbf{Landing pública} con código QR para la captación en eventos presenciales.
\end{itemize}

\paragraph{Interfaces de hardware y comunicación.} El producto no requiere hardware específico. Se
despliega en contenedores sobre un servidor virtual con proxy inverso y certificado TLS. Toda
comunicación externa ocurre sobre HTTPS.

\subsection{Funciones del producto}

\begin{table}[H]
    \centering
    \tablacompacta
    \begin{tabularx}{\textwidth}{p{1.4cm}p{4.0cm}X}
        \toprule
        \cabecera{\textbf{Código}} & \textbf{Función} & \textbf{Resumen} \\
        \midrule
        F-01 & Atención conversacional & Recibir mensajes, mantener el estado del diálogo y responder según el paso actual. \\
        F-02 & Publicación del catálogo & Presentar categorías y programas con costos, duración y modalidad. \\
        F-03 & Asistencia con IA & Responder preguntas abiertas sobre posgrados con base en conocimiento institucional. \\
        F-04 & Captura del prospecto & Solicitar consentimiento y recoger nombre, correo y programa de interés. \\
        F-05 & Entrega al asesor & Abrir la conversación de atención con la ficha del prospecto y mantener el puente. \\
        F-06 & Integración con el CRM & Registrar o actualizar el contacto del prospecto. \\
        F-07 & Registro de captación & Consignar el prospecto en la hoja de cálculo del evento. \\
        F-08 & Administración y métricas & Autenticar asesores, mostrar indicadores y administrar sesiones. \\
        \bottomrule
    \end{tabularx}
    \caption{Funciones del producto.}
\end{table}

\subsection{Características de los usuarios}

\begin{table}[H]
    \centering
    \tablacompacta
    \begin{tabularx}{\textwidth}{p{2.7cm}p{2.4cm}p{2.2cm}X}
        \toprule
        \cabecera{\textbf{Usuario}} & \textbf{Formación} & \textbf{Experiencia técnica} & \textbf{Uso del sistema} \\
        \midrule
        Aspirante a posgrado & Profesional titulado, 24--45 años & Usuario cotidiano de WhatsApp; no requiere formación & Consulta programas, pregunta y deja sus datos. Uso esporádico, desde el móvil, a cualquier hora. \\[0.2cm]
        Asesor de Admisiones & Profesional en mercadeo o áreas afines & Usuario de herramientas de oficina y plataformas de atención & Atiende las conversaciones entregadas por el bot y consulta el tablero. Uso diario en horario laboral. \\[0.2cm]
        Coordinador de Admisiones & Profesional con rol directivo & Usuario intermedio & Consulta indicadores de captación y conversión. Uso semanal. \\[0.2cm]
        Administrador del sistema & Ingeniero de sistemas & Avanzado & Configura, despliega, monitorea y administra cuentas de asesores. Uso según necesidad. \\
        \bottomrule
    \end{tabularx}
    \caption{Características de los usuarios.}
\end{table}

\subsection{Restricciones}

\begin{enumerate}[label=\textbf{R-\arabic*.}, leftmargin=1.5cm, itemsep=3pt]
    \item \textbf{Políticas de la plataforma de mensajería.} El sistema solo puede iniciar
    conversaciones dentro de la ventana de 24 horas posterior al último mensaje del usuario; fuera de
    ella se requieren plantillas aprobadas por Meta.
    \item \textbf{Límites de la interfaz de WhatsApp.} Máximo 10 filas por lista interactiva y 3
    botones por mensaje; el diseño del menú debe respetar esos límites.
    \item \textbf{Marco legal.} El tratamiento de datos personales debe cumplir la Ley 1581 de 2012
    y el Decreto 1377 de 2013: consentimiento previo, expreso e informado, y minimización de datos.
    \item \textbf{Tecnológicas.} Node.js con TypeScript, Express, Redis y despliegue en contenedores
    Docker sobre el servidor virtual provisto.
    \item \textbf{Dependencia de terceros.} La disponibilidad del sistema está acotada por la de
    Meta, el proveedor de IA, Chatwoot y el CRM.
    \item \textbf{Presupuesto y consumo.} El uso del modelo de lenguaje y de la API de mensajería
    tiene costo por uso, lo que obliga a limitar la longitud de las respuestas y a controlar la
    frecuencia de mensajes por usuario.
    \item \textbf{Vigencia de la información.} Los costos de los programas cambian por período
    académico y deben poder actualizarse sin desplegar una nueva versión.
    \item \textbf{Idioma y localización.} Español de Colombia; fechas y horas en la zona horaria
    \textit{America/Bogota}; valores en pesos colombianos.
\end{enumerate}

\subsection{Suposiciones y dependencias}

\begin{itemize}[itemsep=2pt]
    \item Se supone que el aspirante dispone de WhatsApp y conexión a datos.
    \item Se supone que la universidad mantiene el número verificado y la cuenta de WhatsApp
    Business activa.
    \item Se supone que la Coordinación de Admisiones entrega y mantiene actualizada la información
    de costos, duración y requisitos de los programas.
    \item Se supone que existe al menos un asesor disponible en horario laboral para atender los
    prospectos entregados.
    \item El sistema depende de que las credenciales de los servicios externos permanezcan vigentes.
\end{itemize}

\clearpage
\subsection{Requisitos funcionales}

Cada requisito se identifica con el código \textbf{RF-nn} y una prioridad expresada como
\textit{Alta}, \textit{Media} o \textit{Baja}.

{\tablacompacta
\begin{longtable}{p{1.4cm}p{3.3cm}p{7.4cm}p{1.4cm}}
\caption{Requisitos funcionales del flujo de posgrados.}\label{tab:rf}\\
\toprule
\cabecera{\textbf{ID}} & \textbf{Nombre} & \textbf{Descripción} & \textbf{Prior.} \\
\midrule
\endfirsthead
\toprule
\textbf{ID} & \textbf{Nombre} & \textbf{Descripción} & \textbf{Prior.} \\
\midrule
\endhead
\bottomrule
\endfoot
RF-01 & Recepción de mensajes & El sistema debe recibir por webhook los mensajes entrantes de la plataforma de mensajería, verificar el token de validación y responder con acuse inmediato. & Alta \\
RF-02 & Idempotencia de mensajes & El sistema debe procesar cada mensaje una sola vez, descartando reenvíos del mismo identificador durante al menos 5 minutos. & Alta \\
RF-03 & Enrutamiento por estado & El sistema debe entregar cada mensaje al manejador correspondiente al paso actual de la sesión del usuario, y mostrar el menú principal si el paso es desconocido o no existe sesión. & Alta \\
RF-04 & Menú principal & El sistema debe presentar un menú interactivo con las opciones disponibles, incluida la entrada al catálogo de posgrados. & Alta \\
RF-05 & Comando de reinicio & El sistema debe interpretar la palabra \texttt{menu}, en cualquier punto del diálogo, como orden de volver al menú principal. & Media \\
RF-06 & Selección de categoría & El sistema debe permitir elegir entre especializaciones, maestrías y doctorados mediante una lista interactiva. & Alta \\
RF-07 & Detalle del catálogo & El sistema debe mostrar, por categoría, el listado de programas con su valor por semestre, duración, modalidad y valor de inscripción. & Alta \\
RF-08 & Acciones tras el detalle & Tras mostrar un detalle, el sistema debe ofrecer tres acciones: preguntar al asistente de IA, hablar con un asesor o volver a otros programas. & Media \\
RF-09 & Asistente de IA & El sistema debe responder preguntas abiertas sobre posgrados usando un modelo de lenguaje restringido a la base de conocimiento institucional (requisitos, calendario, costos y financiación). & Alta \\
RF-10 & Límite del asistente & Cuando la pregunta exceda la base de conocimiento, el asistente debe declararlo y ofrecer el contacto con un asesor, en lugar de inventar una respuesta. & Alta \\
RF-11 & Consentimiento informado & Antes de capturar cualquier dato personal en el flujo de captación, el sistema debe solicitar autorización explícita de tratamiento de datos e informar su finalidad. & Alta \\
RF-12 & Terminación sin registro & Si el usuario no autoriza el tratamiento de datos, el sistema debe cerrar el diálogo sin persistir ningún dato personal. & Alta \\
RF-13 & Captura del nombre & El sistema debe solicitar el nombre completo y rechazar entradas de menos de 3 caracteres, volviendo a solicitarlo. & Alta \\
RF-14 & Captura del correo & El sistema debe solicitar el correo electrónico, validar su formato y normalizarlo a minúsculas; si no es válido, debe volver a solicitarlo con un ejemplo. & Alta \\
RF-15 & Captura del programa & El sistema debe solicitar el programa o área de interés, admitiendo respuestas específicas o generales. & Alta \\
RF-16 & Persistencia parcial & El sistema debe conservar los datos ya capturados mientras dure la sesión, de modo que el usuario pueda continuar el diálogo sin repetirlos. & Alta \\
RF-17 & Ficha del prospecto & Al entregar el prospecto, el sistema debe componer una ficha con nombre, número de WhatsApp, correo, programa de interés, categoría consultada y fecha de recepción. & Alta \\
RF-18 & Apertura de la atención & El sistema debe crear o recuperar el contacto y abrir una conversación en la plataforma de atención, asignada al equipo de Admisiones, con la etiqueta \texttt{posgrado-prospecto} y prioridad media. & Alta \\
RF-19 & Atributos de la conversación & El sistema debe adjuntar a la conversación los atributos de área, programa de interés y correo, para permitir su búsqueda y segmentación. & Media \\
RF-20 & Conversación puente & Mientras dure la atención humana, el sistema debe reenviar al asesor los mensajes del aspirante y al aspirante las respuestas del asesor, de forma transparente. & Alta \\
RF-21 & Cierre de la atención & Al resolverse la conversación en la plataforma de atención, el sistema debe devolver al usuario al menú principal. & Media \\
RF-22 & Sincronización con el CRM & El sistema debe registrar o actualizar el contacto del prospecto en el CRM, deduplicando por correo, cuando la integración esté activa. & Media \\
RF-23 & Degradación del CRM & Si una propiedad personalizada no existe en el CRM, el sistema debe reintentar el registro solo con las propiedades estándar. & Baja \\
RF-24 & Registro de captación & En el flujo de eventos, el sistema debe agregar una fila a la hoja de cálculo con fecha, WhatsApp, nombre, correo, interés en financiación, programa y autorización. & Media \\
RF-25 & Actualización diferida & El sistema debe poder completar posteriormente las columnas de programa de interés y autorización de la fila ya registrada. & Baja \\
RF-26 & Estado conversacional & El sistema debe almacenar el paso y los datos parciales de cada conversación con un tiempo de vida configurable, tras el cual la sesión expira. & Alta \\
RF-27 & Registro de métricas & El sistema debe contabilizar los eventos relevantes del flujo (consulta de programas, uso del asistente, prospectos completados, entregas a asesor). & Media \\
RF-28 & Exposición de métricas & El sistema debe exponer sus métricas en un punto de acceso compatible con la herramienta de monitoreo, además de un punto de verificación de estado. & Media \\
RF-29 & Autenticación del tablero & El tablero de asesores debe exigir usuario y contraseña, y emitir un token de sesión con vigencia limitada. & Alta \\
RF-30 & Autorización por rol & El tablero debe mostrar la sección de administración de asesores únicamente a los usuarios con rol de administrador. & Media \\
RF-31 & Consulta de indicadores & El tablero debe mostrar los indicadores de uso del bot y el número de prospectos captados por período. & Media \\
RF-32 & Reinicio de sesión & Un administrador debe poder reiniciar la sesión conversacional de un usuario para desbloquear un diálogo. & Baja \\
RF-33 & Landing con código QR & El sistema debe exponer una página pública que, mediante un código QR, inicie la conversación de WhatsApp con el bot. & Baja \\
RF-34 & Activación por configuración & El flujo de captación por evento y la integración con el CRM deben poder activarse o desactivarse por configuración, sin cambios en el código. & Media \\
RF-35 & Manejo de fallos externos & Si un servicio externo falla, el sistema debe informar al usuario en lenguaje comprensible y ofrecer una alternativa de contacto, sin exponer detalles técnicos. & Alta \\
\end{longtable}}

\subsection{Requisitos no funcionales}

{\tablacompacta
\begin{longtable}{p{1.6cm}p{3.1cm}p{7.4cm}p{1.4cm}}
\caption{Requisitos no funcionales.}\label{tab:rnf}\\
\toprule
\cabecera{\textbf{ID}} & \textbf{Atributo} & \textbf{Descripción} & \textbf{Prior.} \\
\midrule
\endfirsthead
\toprule
\textbf{ID} & \textbf{Atributo} & \textbf{Descripción} & \textbf{Prior.} \\
\midrule
\endhead
\bottomrule
\endfoot
RNF-01 & Rendimiento & El acuse de recibo del webhook debe emitirse en menos de 2 segundos para evitar reintentos de la plataforma de mensajería. & Alta \\
RNF-02 & Rendimiento & Las respuestas de catálogo y de navegación deben entregarse al usuario en menos de 3 segundos en el 95\% de los casos. & Alta \\
RNF-03 & Rendimiento & Las respuestas del asistente de IA deben entregarse en menos de 10 segundos; si se supera el umbral, el sistema debe enviar un aviso de espera. & Media \\
RNF-04 & Disponibilidad & El servicio debe estar disponible al menos el 99\% del tiempo mensual, con reinicio automático de los contenedores ante caída. & Alta \\
RNF-05 & Escalabilidad & El sistema debe soportar al menos 200 conversaciones simultáneas sin degradación perceptible, manteniendo el estado fuera del proceso de aplicación. & Media \\
RNF-06 & Seguridad & Toda comunicación externa debe realizarse sobre HTTPS con certificado válido. & Alta \\
RNF-07 & Seguridad & Las credenciales y claves de servicios externos deben residir en variables de entorno y nunca en el código fuente ni en el repositorio. & Alta \\
RNF-08 & Seguridad & El token de sesión del tablero debe transportarse en una cookie de solo servidor, con vigencia máxima de 8 horas. & Alta \\
RNF-09 & Seguridad & El webhook debe rechazar peticiones que no presenten el token de verificación esperado. & Alta \\
RNF-10 & Seguridad & Las contraseñas de los asesores deben almacenarse cifradas mediante función de resumen con sal. & Alta \\
RNF-11 & Privacidad & El sistema debe recoger únicamente los datos necesarios para el contacto y registrar la autorización de tratamiento, conforme a la Ley 1581 de 2012. & Alta \\
RNF-12 & Privacidad & El contenido de los mensajes de los usuarios no debe almacenarse en texto plano ni en registros de aplicación. & Alta \\
RNF-13 & Privacidad & El estado conversacional debe ser efímero, con expiración automática configurable. & Alta \\
RNF-14 & Usabilidad & El usuario debe alcanzar el detalle de un programa en un máximo de tres interacciones desde el mensaje inicial. & Alta \\
RNF-15 & Usabilidad & Los mensajes deben usar lenguaje claro y no técnico, con un máximo aproximado de 1.000 caracteres por mensaje. & Media \\
RNF-16 & Usabilidad & La navegación debe apoyarse en listas y botones interactivos antes que en la escritura libre, para reducir errores de entrada. & Alta \\
RNF-17 & Mantenibilidad & La lógica conversacional debe estar aislada del proveedor de mensajería mediante una interfaz de adaptador, de modo que cambiar de proveedor no afecte los flujos. & Alta \\
RNF-18 & Mantenibilidad & Cada flujo debe residir en un módulo independiente, con un único punto de enrutamiento por paso de sesión. & Media \\
RNF-19 & Portabilidad & El despliegue completo debe reproducirse con contenedores, sin instalación manual de dependencias en el servidor. & Media \\
RNF-20 & Observabilidad & El sistema debe exponer métricas operacionales y de negocio, con tableros y reglas de alerta. & Media \\
RNF-21 & Confiabilidad & El fallo de un servicio externo no debe interrumpir la conversación: el sistema debe degradarse ofreciendo una vía alterna. & Alta \\
RNF-22 & Localización & Fechas, horas y valores monetarios deben presentarse en formato colombiano y zona horaria \textit{America/Bogota}. & Media \\
RNF-23 & Cumplimiento & El sistema debe respetar las políticas de la plataforma de mensajería, incluida la ventana de 24 horas para mensajes libres. & Alta \\
RNF-24 & Costo & El consumo del modelo de lenguaje debe limitarse mediante control de frecuencia por usuario y respuestas acotadas en longitud. & Media \\
\end{longtable}}

\clearpage
% =========================================================
% 5. HISTORIAS DE USUARIO
% =========================================================
\section{Historias de Usuario}

Las historias se redactan en el formato \textit{Como \textless rol\textgreater{} quiero
\textless funcionalidad\textgreater{} para \textless beneficio\textgreater}, con criterios de
aceptación verificables, prioridad MoSCoW (\textit{Debe}, \textit{Debería}, \textit{Podría}) y
estimación en puntos de historia.

\subsection{Épicas}

\begin{table}[H]
    \centering
    \tablacompacta
    \begin{tabularx}{\textwidth}{p{1.5cm}p{4.2cm}X}
        \toprule
        \cabecera{\textbf{ID}} & \textbf{Épica} & \textbf{Descripción} \\
        \midrule
        E-01 & Información de posgrados & Todo lo que permite al aspirante conocer la oferta sin intervención humana. \\
        E-02 & Captación del prospecto & Consentimiento, captura y calificación de los datos del interesado. \\
        E-03 & Atención humana & Entrega del prospecto al asesor y conversación puente. \\
        E-04 & Trazabilidad y gestión & Registro en CRM y hoja de captación, indicadores y administración. \\
        \bottomrule
    \end{tabularx}
    \caption{Épicas del producto.}
\end{table}

\subsection{Catálogo de historias de usuario}

{\tablacompacta
\begin{longtable}{p{1.2cm}p{0.9cm}p{5.2cm}p{5.0cm}p{1.1cm}p{0.7cm}}
\caption{Historias de usuario del flujo de posgrados.}\label{tab:hu}\\
\toprule
\cabecera{\textbf{ID}} & \textbf{Épica} & \textbf{Historia} & \textbf{Criterios de aceptación} & \textbf{Prior.} & \textbf{Pts} \\
\midrule
\endfirsthead
\toprule
\textbf{ID} & \textbf{Épica} & \textbf{Historia} & \textbf{Criterios de aceptación} & \textbf{Prior.} & \textbf{Pts} \\
\midrule
\endhead
\bottomrule
\endfoot
HU-01 & E-01 & Como aspirante quiero escribir al número de la universidad y recibir de inmediato un menú de opciones, para saber qué puedo hacer sin leer instrucciones. & Al enviar cualquier mensaje sin sesión activa se recibe el menú. El menú llega en menos de 3 segundos. Las opciones son seleccionables, no escritas. & Debe & 3 \\
HU-02 & E-01 & Como aspirante quiero elegir entre especializaciones, maestrías y doctorados, para ver solo los programas del nivel que me interesa. & La lista muestra las tres categorías con duración y modalidad. Al elegir una se muestran únicamente sus programas. Una opción permite volver. & Debe & 3 \\
HU-03 & E-01 & Como aspirante quiero ver el costo por semestre y la duración de cada programa, para saber si está a mi alcance antes de contactar a nadie. & Cada programa muestra nombre, valor por semestre y duración. Se indica el valor de inscripción de la categoría. Se distinguen los programas virtuales. & Debe & 5 \\
HU-04 & E-01 & Como aspirante quiero preguntar en mis propias palabras sobre requisitos, fechas o financiación, para resolver dudas que el catálogo no responde. & La respuesta se basa en la base de conocimiento institucional. Si la pregunta excede su alcance, el asistente lo dice y ofrece un asesor. La respuesta llega en menos de 10 segundos. & Debe & 8 \\
HU-05 & E-01 & Como aspirante quiero volver al menú escribiendo una palabra, para no quedar atrapado en un diálogo. & La palabra \texttt{menu} devuelve al menú principal desde cualquier paso. La sesión anterior se descarta. & Debería & 2 \\
HU-06 & E-02 & Como aspirante quiero que me pidan autorización antes de tomar mis datos, para saber qué van a hacer con ellos. & Antes de la primera pregunta personal se solicita autorización e informa la finalidad. Si se niega, el diálogo termina y no se registra nada. & Debe & 3 \\
HU-07 & E-02 & Como aspirante quiero que el bot me avise si escribo mal mi correo, para no perder el contacto por un error de digitación. & Un correo con formato inválido se rechaza con un ejemplo y se vuelve a solicitar. El correo se guarda en minúsculas. Un nombre de menos de 3 caracteres se rechaza. & Debe & 3 \\
HU-08 & E-02 & Como aspirante quiero poder decir que aún no sé qué programa me interesa, para no quedarme fuera del proceso por indecisión. & El campo de programa admite texto libre, específico o general. No hay validación que bloquee respuestas como ``no sé aún''. & Debería & 1 \\
HU-09 & E-02 & Como aspirante quiero retomar la conversación donde la dejé, para no repetir mis datos si me interrumpen. & Los datos capturados persisten durante la vigencia de la sesión. Al expirar la sesión se vuelve al menú. & Debería & 5 \\
HU-10 & E-03 & Como aspirante quiero hablar con una persona cuando lo necesite, para resolver lo que el bot no puede. & La opción está disponible tras el detalle de un programa y en el menú. Se confirma la conexión con un mensaje. & Debe & 5 \\
HU-11 & E-03 & Como aspirante quiero seguir escribiendo por WhatsApp mientras me atiende un asesor, para no cambiar de aplicación. & Los mensajes del aspirante llegan al asesor y las respuestas del asesor llegan a WhatsApp. Al cerrarse la atención se vuelve al menú. & Debe & 8 \\
HU-12 & E-03 & Como asesor de Admisiones quiero recibir la ficha del prospecto al abrirse la conversación, para no preguntar de nuevo lo que ya respondió. & La conversación inicia con nombre, WhatsApp, correo, programa de interés, categoría vista y fecha. Se asigna al equipo de Admisiones. & Debe & 5 \\
HU-13 & E-03 & Como asesor quiero que las conversaciones de posgrado lleguen etiquetadas y priorizadas, para distinguirlas del resto de la bandeja. & La conversación se crea con la etiqueta \texttt{posgrado-prospecto}, prioridad media y atributos de área, programa y correo. & Debería & 3 \\
HU-14 & E-04 & Como coordinador de Admisiones quiero que cada prospecto quede en el CRM, para hacer seguimiento comercial sin transcribir datos. & El contacto se crea o actualiza deduplicando por correo. Si una propiedad personalizada falta, se reintenta con las estándar. La integración es activable por configuración. & Debería & 5 \\
HU-15 & E-04 & Como coordinador quiero un registro tabular de los prospectos de un evento, para revisarlos sin entrar al sistema. & Cada registro agrega una fila con fecha, WhatsApp, nombre, correo, interés en financiación, programa y autorización. & Debería & 3 \\
HU-16 & E-04 & Como coordinador quiero ver cuántas personas consultaron programas y cuántas dejaron sus datos, para medir el embudo de captación. & El tablero muestra los contadores del período. Los eventos de consulta, uso del asistente y prospectos completados se registran. & Debería & 5 \\
HU-17 & E-04 & Como asesor quiero entrar al tablero con usuario y contraseña, para acceder solo a lo que me corresponde. & El acceso exige credenciales válidas. La sesión caduca a las 8 horas. La sección de asesores solo la ve el administrador. & Debe & 5 \\
HU-18 & E-04 & Como administrador quiero reiniciar la sesión de un usuario, para desbloquear a alguien que quedó atrapado en un paso. & El administrador localiza al usuario por número y reinicia su sesión. El usuario recibe el menú en su siguiente mensaje. & Podría & 2 \\
HU-19 & E-04 & Como coordinador quiero un código QR que abra la conversación con el bot, para usarlo en piezas impresas de los eventos. & Existe una página pública con el código. Al escanearlo se abre WhatsApp con el número y un mensaje inicial. & Podría & 2 \\
HU-20 & E-04 & Como administrador quiero que el bot no se caiga cuando falle un servicio externo, para no perder aspirantes por un problema técnico. & Ante un fallo externo el usuario recibe un mensaje comprensible y una vía alterna de contacto. No se muestran detalles técnicos. El error queda registrado. & Debe & 5 \\
\end{longtable}}

\nota{\textbf{Total estimado:} 81 puntos de historia. Las historias marcadas como \textit{Debe}
(46 puntos) constituyen el producto mínimo viable del flujo de posgrados.}

\clearpage

% =========================================================
% 6. CASOS DE USO
% =========================================================
\section{Diagrama de Casos de Uso}

\subsection{Actores}

\begin{table}[H]
    \centering
    \tablacompacta
    \begin{tabularx}{\textwidth}{p{3.4cm}p{2.0cm}X}
        \toprule
        \cabecera{\textbf{Actor}} & \textbf{Tipo} & \textbf{Descripción} \\
        \midrule
        Aspirante & Primario & Persona interesada en un programa de posgrado que interactúa por WhatsApp. \\
        Asesor de Admisiones & Primario & Atiende las conversaciones entregadas por el bot desde la plataforma de atención. \\
        Coordinador de Admisiones & Primario & Consulta indicadores y el registro de prospectos captados. \\
        Administrador del sistema & Primario & Configura, monitorea y administra cuentas y sesiones. \\
        Plataforma de mensajería & Secundario & Entrega los mensajes entrantes y transporta los salientes. \\
        Servicio de IA & Secundario & Genera las respuestas del asistente. \\
        Plataforma de atención & Secundario & Aloja los contactos y las conversaciones con asesores. \\
        CRM institucional & Secundario & Recibe el contacto del prospecto. \\
        Servicio de hojas de cálculo & Secundario & Recibe el registro tabular de la captación. \\
        \bottomrule
    \end{tabularx}
    \caption{Actores del sistema.}
\end{table}

\subsection{Diagrama}

\begin{figure}[H]
\centering
\resizebox{0.97\textwidth}{!}{%
\begin{tikzpicture}

% --- Límite del sistema ---
\draw[draw=utbazul, thick, rounded corners=4pt] (-1.7,1.0) rectangle (8.7,-9.9);
\node[font=\scriptsize\bfseries, text=utbazul] at (3.5,0.5) {Tooli: flujo de posgrados};

% --- Casos de uso: columna izquierda (aspirante) ---
\node[uc] (u1) at (0.4,-0.9) {UC-01\\ Consultar catálogo\\ de posgrados};
\node[uc] (u2) at (0.4,-2.2) {UC-02\\ Ver detalle de\\ un programa};
\node[uc] (u3) at (0.4,-3.5) {UC-03\\ Preguntar al\\ asistente de IA};
\node[uc] (u4) at (0.4,-4.8) {UC-04\\ Solicitar\\ asesor humano};
\node[uc] (u5) at (0.4,-6.1) {UC-05\\ Registrar\\ prospecto};
\node[uc] (u6) at (0.4,-7.4) {UC-06\\ Otorgar\\ consentimiento};
\node[uc] (u7) at (0.4,-8.7) {UC-07\\ Conversar\\ con el asesor};

% --- Casos de uso: columna derecha (personal y sistemas) ---
\node[uc] (u9)  at (6.6,-0.9) {UC-09\\ Autenticarse\\ en el tablero};
\node[uc] (u10) at (6.6,-2.2) {UC-10\\ Consultar\\ indicadores};
\node[uc] (u11) at (6.6,-3.5) {UC-11\\ Administrar\\ asesores};
\node[uc] (u12) at (6.6,-4.8) {UC-12\\ Reiniciar sesión\\ de un usuario};
\node[uc] (u13) at (6.6,-6.1) {UC-13\\ Sincronizar\\ con el CRM};
\node[uc] (u14) at (6.6,-7.4) {UC-14\\ Registrar en hoja\\ de captación};
\node[uc] (u8)  at (6.6,-8.7) {UC-08\\ Atender al\\ prospecto};

% --- Actores ---
\actor{-3.7}{-4.0}{Aspirante}
\actor{10.7}{-1.6}{Coordinador de\\ Admisiones}
\actor{10.7}{-4.4}{Administrador\\ del sistema}
\actor{10.7}{-8.0}{Asesor de\\ Admisiones}

% --- Asociaciones del aspirante ---
\draw[draw=grismedio] (-3.35,-4.0) -- (u1.west);
\draw[draw=grismedio] (-3.35,-4.0) -- (u2.west);
\draw[draw=grismedio] (-3.35,-4.0) -- (u3.west);
\draw[draw=grismedio] (-3.35,-4.0) -- (u4.west);
\draw[draw=grismedio] (-3.35,-4.0) -- (u6.west);
\draw[draw=grismedio] (-3.35,-4.0) -- (u7.west);

% --- Asociaciones del personal ---
\draw[draw=grismedio] (10.35,-1.6) -- (u9.east);
\draw[draw=grismedio] (10.35,-1.6) -- (u10.east);
\draw[draw=grismedio] (10.35,-4.4) -- (u9.east);
\draw[draw=grismedio] (10.35,-4.4) -- (u11.east);
\draw[draw=grismedio] (10.35,-4.4) -- (u12.east);
\draw[draw=grismedio] (10.35,-8.0) -- (u8.east);
\draw[draw=grismedio] (10.35,-8.0) -- (u9.east);

% --- Relaciones include / extend ---
\draw[flechad] (u4.south) -- node[font=\tiny, right=1pt] {$\ll$include$\gg$} (u5.north);
\draw[flechad] (u5.south) -- node[font=\tiny, right=1pt] {$\ll$include$\gg$} (u6.north);
\draw[flechad] (u3.north) -- node[font=\tiny, right=1pt] {$\ll$extend$\gg$} (u2.south);
\draw[flechad] (u5.east) -- node[font=\tiny, above, pos=0.5] {$\ll$include$\gg$} (u13.west);
\draw[flechad] (u5.east) -- node[font=\tiny, below, pos=0.35] {$\ll$extend$\gg$} (u14.west);
\draw[flechad] (u7.east) -- node[font=\tiny, above, pos=0.5] {$\ll$include$\gg$} (u8.west);

\end{tikzpicture}}
\caption{Diagrama de casos de uso del flujo de posgrados.}
\label{fig:casos-uso}
\end{figure}

\subsection{Especificación de los casos de uso principales}

\subsubsection{UC-01. Consultar catálogo de posgrados}

\begin{table}[H]
    \centering
    \tablacompacta
    \begin{tabularx}{\textwidth}{p{3.2cm}X}
        \toprule
        \cabecera{\textbf{Campo}} & \textbf{Contenido} \\
        \midrule
        \textbf{Actor principal} & Aspirante. \\
        \textbf{Actores secundarios} & Plataforma de mensajería. \\
        \textbf{Descripción} & El aspirante consulta la oferta de posgrados por categoría. \\
        \textbf{Precondiciones} & El aspirante ha iniciado conversación con el bot. \\
        \textbf{Postcondiciones} & Se registra el evento de consulta y la categoría queda guardada en la sesión. \\
        \textbf{Flujo principal} &
        1. El aspirante selecciona ``Ver programas'' en el menú. \newline
        2. El sistema muestra la lista de categorías: especializaciones, maestrías y doctorados. \newline
        3. El aspirante selecciona una categoría. \newline
        4. El sistema muestra los programas de esa categoría con valor por semestre, duración y modalidad. \newline
        5. El sistema ofrece las acciones de continuación. \\
        \textbf{Flujos alternativos} &
        3a. El aspirante escribe \texttt{menu}: el sistema regresa al menú principal y termina el caso. \newline
        3b. El aspirante envía una opción no válida: el sistema repite la lista de categorías. \\
        \textbf{Excepciones} & E1. Falla el envío del mensaje: el sistema reintenta una vez y registra el error. \\
        \textbf{Requisitos} & RF-06, RF-07, RF-08, RF-27. \\
        \bottomrule
    \end{tabularx}
    \caption{Especificación del caso de uso UC-01.}
\end{table}

\subsubsection{UC-03. Preguntar al asistente de IA}

\begin{table}[H]
    \centering
    \tablacompacta
    \begin{tabularx}{\textwidth}{p{3.2cm}X}
        \toprule
        \cabecera{\textbf{Campo}} & \textbf{Contenido} \\
        \midrule
        \textbf{Actor principal} & Aspirante. \\
        \textbf{Actores secundarios} & Servicio de IA. \\
        \textbf{Descripción} & El aspirante formula preguntas abiertas y recibe respuestas basadas en la base de conocimiento institucional. \\
        \textbf{Precondiciones} & El aspirante seleccionó la opción de asistente de IA. \\
        \textbf{Postcondiciones} & La sesión queda en el paso de conversación con IA; se registra el evento de uso. \\
        \textbf{Flujo principal} &
        1. El sistema informa que el asistente está activo y cómo salir. \newline
        2. El aspirante escribe su pregunta. \newline
        3. El sistema consulta el modelo de lenguaje con el contexto de la base de conocimiento. \newline
        4. El sistema entrega la respuesta y queda a la espera de nuevas preguntas. \\
        \textbf{Flujos alternativos} &
        2a. La pregunta excede la base de conocimiento: el sistema lo declara y ofrece contacto con un asesor. \newline
        2b. El aspirante escribe \texttt{menu}: el sistema cierra el asistente y vuelve al menú. \\
        \textbf{Excepciones} &
        E1. El servicio de IA no responde o excede el tiempo límite: el sistema informa la demora y ofrece hablar con un asesor. \newline
        E2. Se supera el límite de mensajes por usuario: el sistema pide esperar unos minutos. \\
        \textbf{Requisitos} & RF-09, RF-10, RF-35, RNF-03, RNF-24. \\
        \bottomrule
    \end{tabularx}
    \caption{Especificación del caso de uso UC-03.}
\end{table}

\subsubsection{UC-05. Registrar prospecto}

\begin{table}[H]
    \centering
    \tablacompacta
    \begin{tabularx}{\textwidth}{p{3.2cm}X}
        \toprule
        \cabecera{\textbf{Campo}} & \textbf{Contenido} \\
        \midrule
        \textbf{Actor principal} & Aspirante. \\
        \textbf{Actores secundarios} & Plataforma de atención, CRM institucional, servicio de hojas de cálculo. \\
        \textbf{Descripción} & El sistema captura los datos de contacto del aspirante interesado y los deja disponibles para la asesoría y el seguimiento. \\
        \textbf{Precondiciones} & El aspirante solicitó hablar con un asesor (UC-04) y otorgó su consentimiento (UC-06). \\
        \textbf{Postcondiciones} & Existe un prospecto con nombre, correo y programa de interés; el contacto queda sincronizado con el CRM. \\
        \textbf{Flujo principal} &
        1. El sistema solicita el nombre completo. \newline
        2. El aspirante responde; el sistema valida la longitud mínima y lo guarda en la sesión. \newline
        3. El sistema solicita el correo electrónico. \newline
        4. El aspirante responde; el sistema valida el formato, lo normaliza y lo guarda. \newline
        5. El sistema solicita el programa o área de interés. \newline
        6. El aspirante responde; el sistema arma la ficha del prospecto. \newline
        7. El sistema sincroniza el contacto con el CRM (UC-13). \newline
        8. El sistema entrega el prospecto al asesor (UC-08) y confirma al aspirante. \\
        \textbf{Flujos alternativos} &
        2a. Nombre de menos de 3 caracteres: el sistema lo rechaza y vuelve a solicitarlo. \newline
        4a. Correo con formato inválido: el sistema lo rechaza mostrando un ejemplo y vuelve a solicitarlo. \newline
        6a. El aspirante responde de forma general (``no sé aún''): el sistema acepta la respuesta. \newline
        7a. La integración con el CRM está desactivada: el paso se omite sin afectar el flujo. \\
        \textbf{Excepciones} &
        E1. Falla la creación de la conversación de atención: el sistema informa al aspirante, entrega un canal alterno de contacto y registra el error. \newline
        E2. La sesión expira durante la captura: los datos parciales se pierden y el aspirante vuelve al menú. \\
        \textbf{Requisitos} & RF-11 a RF-19, RF-22, RF-24, RF-26, RF-35. \\
        \bottomrule
    \end{tabularx}
    \caption{Especificación del caso de uso UC-05.}
\end{table}

\subsubsection{UC-07. Conversar con el asesor}

\begin{table}[H]
    \centering
    \tablacompacta
    \begin{tabularx}{\textwidth}{p{3.2cm}X}
        \toprule
        \cabecera{\textbf{Campo}} & \textbf{Contenido} \\
        \midrule
        \textbf{Actor principal} & Aspirante. \\
        \textbf{Actores secundarios} & Asesor de Admisiones, plataforma de atención. \\
        \textbf{Descripción} & El aspirante y el asesor conversan a través del bot, que actúa como puente entre WhatsApp y la plataforma de atención. \\
        \textbf{Precondiciones} & Existe una conversación abierta asociada a la sesión del aspirante. \\
        \textbf{Postcondiciones} & La conversación queda resuelta y la sesión regresa al menú principal. \\
        \textbf{Flujo principal} &
        1. El aspirante escribe un mensaje. \newline
        2. El sistema lo publica en la conversación de la plataforma de atención. \newline
        3. El asesor responde desde la plataforma. \newline
        4. El sistema recibe la respuesta por webhook y la envía al aspirante por WhatsApp. \newline
        5. Los pasos 1 a 4 se repiten hasta que el asesor cierra la conversación. \newline
        6. El sistema devuelve al aspirante al menú principal. \\
        \textbf{Flujos alternativos} &
        1a. El aspirante escribe \texttt{menu}: el sistema cierra el puente y regresa al menú. \newline
        3a. La conversación ocurre fuera del horario de atención: el sistema informa el horario y deja el mensaje en cola. \\
        \textbf{Excepciones} & E1. La plataforma de atención no está disponible: el sistema informa al aspirante y ofrece un canal alterno. \\
        \textbf{Requisitos} & RF-20, RF-21, RF-35, RNF-21. \\
        \bottomrule
    \end{tabularx}
    \caption{Especificación del caso de uso UC-07.}
\end{table}

\clearpage
% =========================================================
% 7. MOCKUP
% =========================================================
\section{Mockup de la Aplicación}

La interfaz principal del producto es \textbf{conversacional}: no tiene pantallas propias, sino
mensajes dentro de WhatsApp. Por eso el mockup se presenta como una secuencia de capturas del
recorrido real del aspirante. A ellas se suman las pantallas del \textbf{tablero web de asesores} y
de la \textbf{bandeja de atención humana}, que sí son interfaces gráficas convencionales.

Las imágenes de esta sección provienen del prototipo funcional ya desplegado; por tanto, el mockup
no es una aproximación de diseño, sino la interfaz construida y contrastada con los requisitos
especificados en la sección 4.

\subsection{Pantallas conversacionales (WhatsApp)}

\begin{figure}[H]
\centering
\begin{minipage}{0.46\linewidth}
  \centering
  \captura[height=10.2cm]{whatsapp-registro-datos.png}\\[2mm]
  {\small (a) Apertura del flujo y captura de datos.}
\end{minipage}\hfill
\begin{minipage}{0.46\linewidth}
  \centering
  \captura[height=10.2cm]{whatsapp-registro-confirmacion.png}\\[2mm]
  {\small (b) Programa de interés y confirmación.}
\end{minipage}
\caption{Interfaz conversacional de Tooli en WhatsApp: registro del prospecto de posgrados.}
\label{fig:mockup-whatsapp}
\end{figure}

\paragraph{(a) Apertura del flujo y captura de datos.} El bot recibe al aspirante con el menú de
posgrados y, al seleccionar \textit{Registrarme}, abre el formulario conversacional: primero informa
la finalidad del registro ---enviar información sobre programas, becas y novedades---, luego solicita
el \textbf{nombre completo} y, tras confirmarlo, el \textbf{correo electrónico}. Cada mensaje recuerda
que la palabra \texttt{menu} cancela el diálogo, de modo que el usuario nunca queda atrapado en un
paso. La pantalla materializa los requisitos RF-04, RF-05, RF-11, RF-13 y RF-14, y las historias
HU-01, HU-05, HU-06 y HU-07.

\paragraph{(b) Programa de interés y confirmación.} La elección del tipo de posgrado se resuelve con
una \textbf{lista interactiva} (\textit{Maestría $\cdot$ 3--4 semestres $\cdot$ presencial y
virtual}) en lugar de escritura libre, lo que elimina los errores de entrada. Al cerrar el flujo, el
bot devuelve la \textbf{ficha del prospecto} ---nombre, correo, número de WhatsApp y programa--- que
es exactamente la información que recibirá el asesor, y anuncia el contacto posterior. Corresponde a
los requisitos RF-06, RF-15, RF-16 y RF-17, a los no funcionales RNF-14 y RNF-16, y a las historias
HU-02, HU-08 y HU-12.

\subsection{Tablero web de asesores}

\subsubsection{Inicio de sesión}

\begin{figure}[H]
\centering
\captura[width=0.92\linewidth]{panel-login.png}
\caption{Pantalla de acceso al tablero de asesores.}
\label{fig:mockup-login}
\end{figure}

El tablero es de uso interno: no admite registro público y exige usuario y contraseña de una cuenta
creada por el administrador. La sesión se emite como token con vigencia limitada y las contraseñas se
almacenan cifradas, nunca en claro. Implementa el requisito RF-29 y los no funcionales RNF-08 y
RNF-10, correspondientes a la historia HU-17.

\subsubsection{Panel de inicio}

\begin{figure}[H]
\centering
\captura[width=\linewidth]{panel-inicio.png}
\caption{Panel de inicio: indicadores del día, alertas y accesos rápidos.}
\label{fig:mockup-inicio}
\end{figure}

Es la vista de entrada del asesor. La fila superior resume los indicadores del día ---mensajes
atendidos, preguntas resueltas por la IA, prospectos captados y casos escalados a un asesor---, que
son los contadores del embudo descritos en la sección 4. El bloque \textit{Requiere atención} señala
las conversaciones sin asignar en la plataforma de atención, y la columna de accesos rápidos enlaza
los módulos operativos (conversaciones, consulta de turno, recibo, CRM, posgrados, recordatorios y
métricas). Cubre los requisitos RF-27, RF-28 y RF-31, y la historia HU-16.

\subsubsection{Gestión de usuarios}

\begin{figure}[H]
\centering
\captura[width=\linewidth]{panel-usuarios.png}
\caption{Gestión de usuarios: administración de las cuentas de asesor.}
\label{fig:mockup-usuarios}
\end{figure}

Esta sección solo es visible para el rol de \textbf{administrador}. Desde ella se crean cuentas de
asesor, se asigna su área de trabajo y se consulta su actividad (acciones realizadas y casos a
cargo), lo que permite revocar accesos sin eliminar el registro histórico. Implementa los requisitos
RF-30 y RF-32, y las historias HU-17 y HU-18.

\subsection{Bandeja de atención humana}

\begin{figure}[H]
\centering
\captura[width=\linewidth]{chatwoot-bandeja.png}
\caption{Bandeja de la plataforma de atención: destino del prospecto entregado por el bot.}
\label{fig:mockup-chatwoot}
\end{figure}

Es el punto donde aterriza la entrega del prospecto (\textit{handoff}). Las conversaciones llegan al
buzón de WhatsApp del equipo y se distribuyen entre las bandejas \textit{Mine}, \textit{Unassigned} y
\textit{All}, con equipos diferenciados (admisiones y soporte técnico) y etiquetas de clasificación.
El asesor responde desde aquí y el bot reenvía su respuesta al aspirante por WhatsApp sin que este
cambie de aplicación. Corresponde a los requisitos RF-18 a RF-21 y a las historias HU-11, HU-12 y
HU-13.

\subsection{Trazabilidad de las pantallas}

\begin{table}[H]
    \centering
    \tablacompacta
    \begin{tabularx}{\textwidth}{p{4.6cm}Xp{2.6cm}}
        \toprule
        \cabecera{\textbf{Pantalla}} & \textbf{Requisitos} & \textbf{Historias} \\
        \midrule
        Registro en WhatsApp (figura \ref{fig:mockup-whatsapp}a) & RF-04, RF-05, RF-11, RF-13, RF-14 & HU-01, HU-05, HU-06, HU-07 \\
        Confirmación del prospecto (figura \ref{fig:mockup-whatsapp}b) & RF-06, RF-15, RF-16, RF-17, RNF-14, RNF-16 & HU-02, HU-08, HU-12 \\
        Inicio de sesión (figura \ref{fig:mockup-login}) & RF-29, RNF-08, RNF-10 & HU-17 \\
        Panel de inicio (figura \ref{fig:mockup-inicio}) & RF-27, RF-28, RF-31 & HU-16 \\
        Gestión de usuarios (figura \ref{fig:mockup-usuarios}) & RF-30, RF-32 & HU-17, HU-18 \\
        Bandeja de atención (figura \ref{fig:mockup-chatwoot}) & RF-18, RF-19, RF-20, RF-21 & HU-11, HU-12, HU-13 \\
        \bottomrule
    \end{tabularx}
    \caption{Correspondencia entre las pantallas del prototipo y la especificación.}
\end{table}

\nota{\textbf{Nota de diseño.} Las pantallas conversacionales respetan las restricciones reales de la
plataforma de mensajería: máximo 10 filas por lista interactiva y 3 botones por mensaje. Por eso la
selección de programa se entrega como lista desplegable (\textit{Ver opciones}) y la navegación se
apoya en listas y botones, nunca en escritura libre, salvo en la captura de datos y en el asistente
de IA.}

\clearpage
% =========================================================
% 8. TRAZABILIDAD
% =========================================================
\section{Matriz de Trazabilidad}

La matriz relaciona cada historia de usuario con los requisitos funcionales que la implementan y con
los casos de uso que la modelan. Permite verificar que ningún requisito quede huérfano y que ninguna
necesidad quede sin cubrir.

{\tablacompacta
\begin{longtable}{p{1.5cm}p{5.2cm}p{5.0cm}p{3.0cm}}
\caption{Trazabilidad entre historias de usuario, requisitos y casos de uso.}\label{tab:trazabilidad}\\
\toprule
\cabecera{\textbf{HU}} & \textbf{Necesidad} & \textbf{Requisitos funcionales} & \textbf{Casos de uso} \\
\midrule
\endfirsthead
\toprule
\textbf{HU} & \textbf{Necesidad} & \textbf{Requisitos funcionales} & \textbf{Casos de uso} \\
\midrule
\endhead
\bottomrule
\endfoot
HU-01 & Menú inmediato al escribir & RF-01, RF-02, RF-03, RF-04 & UC-01 \\
HU-02 & Selección de categoría & RF-06 & UC-01 \\
HU-03 & Costos y duración visibles & RF-07, RF-08 & UC-01, UC-02 \\
HU-04 & Preguntas abiertas resueltas & RF-09, RF-10 & UC-03 \\
HU-05 & Reinicio del diálogo & RF-05 & UC-01, UC-03, UC-07 \\
HU-06 & Consentimiento previo & RF-11, RF-12 & UC-06 \\
HU-07 & Validación de datos & RF-13, RF-14 & UC-05 \\
HU-08 & Respuesta general admitida & RF-15 & UC-05 \\
HU-09 & Continuidad del diálogo & RF-16, RF-26 & UC-05 \\
HU-10 & Acceso a un asesor & RF-08, RF-17 & UC-04 \\
HU-11 & Conversación puente & RF-20, RF-21 & UC-07 \\
HU-12 & Ficha del prospecto & RF-17, RF-18 & UC-05, UC-08 \\
HU-13 & Etiqueta y prioridad & RF-18, RF-19 & UC-08 \\
HU-14 & Prospecto en el CRM & RF-22, RF-23, RF-34 & UC-13 \\
HU-15 & Registro tabular del evento & RF-24, RF-25, RF-34 & UC-14 \\
HU-16 & Indicadores del embudo & RF-27, RF-28, RF-31 & UC-10 \\
HU-17 & Acceso autenticado & RF-29, RF-30 & UC-09, UC-11 \\
HU-18 & Reinicio de sesión ajena & RF-32 & UC-12 \\
HU-19 & Código QR de captación & RF-33 & UC-01 \\
HU-20 & Tolerancia a fallos externos & RF-35 & UC-03, UC-05, UC-07 \\
\end{longtable}}

% =========================================================
% 9. GLOSARIO
% =========================================================
\section{Glosario}

\begin{table}[H]
    \centering
    \tablacompacta
    \begin{tabularx}{\textwidth}{p{3.6cm}X}
        \toprule
        \cabecera{\textbf{Término}} & \textbf{Definición} \\
        \midrule
        Prospecto & Persona interesada que ya entregó sus datos de contacto y quedó registrada para seguimiento. \\
        Handoff & Transferencia de la conversación desde el bot hacia un asesor humano, conservando el contexto. \\
        Máquina de estados & Modelo en el que la conversación avanza por pasos; cada paso determina qué manejador atiende el siguiente mensaje. \\
        Sesión & Estado temporal de la conversación de un usuario, con tiempo de vida limitado. \\
        Webhook & Punto de acceso HTTP que un servicio externo invoca para notificar un evento. \\
        Base de conocimiento & Conjunto de textos institucionales que acota las respuestas del asistente de IA. \\
        Ventana de 24 horas & Período tras el último mensaje del usuario durante el cual la plataforma permite enviar mensajes libres. \\
        Puntos de historia & Unidad relativa de esfuerzo usada para estimar historias de usuario. \\
        MoSCoW & Técnica de priorización: \textit{Must}, \textit{Should}, \textit{Could}, \textit{Won't}. \\
        \bottomrule
    \end{tabularx}
    \caption{Glosario de términos.}
\end{table}

% =========================================================
% 10. CONCLUSIONES
% =========================================================
\section{Conclusiones}

\begin{itemize}[itemsep=4pt]
    \item El análisis del flujo de posgrados de Tooli identificó \textbf{nueve elementos de
    información} centrales, de los cuales solo tres persisten datos personales (prospecto, contacto
    en el CRM y registro de captación). El resto es información volátil o de catálogo, lo que reduce
    de forma deliberada la superficie de datos sensibles del sistema.
    \item El mapa de procesos muestra que el sistema \textbf{no reemplaza el proceso de captación}:
    automatiza por completo sus tres primeras etapas (atracción, información y calificación) y
    habilita la cuarta entregando al asesor humano un prospecto con contexto. El valor del producto
    está tanto en lo que automatiza como en lo que deja de hacer.
    \item La especificación bajo IEEE 830 arrojó \textbf{35 requisitos funcionales y 24 no
    funcionales}. Es significativo que los requisitos no funcionales de privacidad, seguridad y
    tolerancia a fallos concentren la mayor prioridad: en un canal público como WhatsApp, el
    cumplimiento legal y la degradación elegante ante fallos son tan críticos como la funcionalidad.
    \item Las 20 historias de usuario suman \textbf{81 puntos de historia}, de los cuales 46
    corresponden a historias obligatorias. Ese subconjunto define un producto mínimo viable
    verificable: informar el catálogo, resolver dudas con IA y entregar el prospecto al asesor.
    \item El diagrama de casos de uso evidencia que el sistema tiene \textbf{un actor primario
    dominante} (el aspirante, asociado a seis de los catorce casos) y que las relaciones de
    inclusión concentran el riesgo en UC-05, el caso de uso de registro del prospecto, que integra
    tres servicios externos.
    \item El mockup confirma que las restricciones de la plataforma de mensajería (10 filas por
    lista, 3 botones por mensaje, ventana de 24 horas) son un condicionante de diseño y no un detalle
    de implementación: definen la forma del menú, la extensión de los mensajes y el momento en que el
    sistema puede iniciar una conversación.
\end{itemize}

% =========================================================
% 11. REFERENCIAS
% =========================================================
\section{Referencias}

\begin{enumerate}[itemsep=3pt]
    \item IEEE. \textit{IEEE Recommended Practice for Software Requirements Specifications
    (IEEE Std 830-1998)}. Institute of Electrical and Electronics Engineers, 1998.
    \item Congreso de la República de Colombia. \textit{Ley 1581 de 2012, por la cual se dictan
    disposiciones generales para la protección de datos personales}. Bogotá, 2012.
    \item Presidencia de la República de Colombia. \textit{Decreto 1377 de 2013, por el cual se
    reglamenta parcialmente la Ley 1581 de 2012}. Bogotá, 2013.
    \item Object Management Group. \textit{Unified Modeling Language (UML) Specification},
    versión 2.5.1, 2017.
    \item Cohn, M. \textit{User Stories Applied: For Agile Software Development}.
    Addison-Wesley, 2004.
    \item Meta Platforms. \textit{WhatsApp Business Platform: Cloud API Documentation}.
    \url{https://developers.facebook.com/docs/whatsapp/cloud-api}
    \item Universidad Tecnológica de Bolívar. \textit{Portafolio de programas de posgrado}.
    \url{https://www.utb.edu.co}
\end{enumerate}

\end{document}
